\subsection{Реализация исполнителя}

В общем виде процесс исполнения тесно связан с этапом семантического анализа.
Выполняется рекурсивный обход абстрактного синтаксического дерева.
Первым шагом каждое выражение проходит семантическую проверку.
После успешного завершения семантического анализа выражения из AST передаются на этап их вычисления.
Разнотипные выражения обрабатываются по-разному, однако результат вычисления всегда представляет собой некоторый тип данных, представленный в виде объекта -- внутреннего представления.

Окружение для хранения информации о переменных в программном коде реализовано в виде структуры,
содержащей поля с хэш-картой с самими переменными и их значениями, а также ссылка на эту же структуру для организации области видимости при вызове функций.
На рисунке~\ref{f:code_envStruct} приведен пример программной реализации данной структуры.

\begin{figure}[ht]
	\centering
	\vspace{\toppaddingoffigure}
	\begin{lstlisting}[
        language=Go
    ]
type Env struct {
    store map[string]Object
    outer *Env
}
\end{lstlisting}
	\caption{Реализация окружения}
	\label{f:code_envStruct}
\end{figure}

За обработку узлов AST отвечает единственная функция, которая определяет тип узла и передает управление соответствующей функции,
которая после прохождения семантической проверки вычисляет значение для узла данного типа.
В конечном итоге формируется внутреннее представление в виде объектов с примитивными типами данных, либо объекты представляющие составные типы, состоящие из примитивных, так как массивы и хеш-карты.
Фрагмент кода основной функции получения и обработки узлов AST приведен на рисунке~\ref{f:code_evalFragment}.

При обнаружении в коде определения переменной, ее необходимо сохранить в памяти, чтобы в дальнейшем иметь к ней доступ.
Для этого используется окружение.
Переменные в окружении хранятся в виде хэш-карты, ключи которой представляют идентификатор переменной, а значения – внутреннее представление значений переменной, то есть объекты.

Код функций записи переменной в окружение и получения из него приведен на рисунке~\ref{f:code_getsetEnv}.

\clearpage

\begin{figure}[!htb]
	\centering
	\begin{lstlisting}[
        language=Go,
        xleftmargin=.08\textwidth,
        xrightmargin=.08\textwidth
    ]
func Eval(n ast.Node, env *object.Env) object.Object {
    switch node := n.(type) {
    case *ast.Program:
        return evalProgram(node, env)
    case *ast.BlockStatement:
        return evalBlockStatement(node, env)
    case *ast.ExpressionStatement:
        return Eval(node.Expression, env)
    case *ast.ReturnStatement:
        val := Eval(node.Value, env)
        if isError(val) {
            return val
        }

        return &object.Return{Value: val}
    case *ast.AssignStatement:
        val := Eval(node.Value, env)
        if isError(val) {
            return val
        }

        env.Set(node.Name.Value, val)
    case *ast.IntegerLiteral:
        return &object.Integer{Value: node.Value}
    case *ast.Boolean:
        if node.Value {
            return TRUE
        }
        return FALSE
\end{lstlisting}
	\caption{Фрагмент кода исполнителя}
	\label{f:code_evalFragment}
\end{figure}

\begin{figure}[!htb]
	\centering
	\begin{lstlisting}[
        language=Go,
        xleftmargin=.08\textwidth,
        xrightmargin=.08\textwidth
    ]
func (e *Env) Get(key string) (Object, bool) {
    obj, ok := e.store[key]
    if !ok && e.outer != nil {
        obj, ok = e.outer.Get(key)
    }
    return obj, ok
}
func (e *Env) Set(key string, val Object) Object {
    e.store[key] = val
    return val
}    
\end{lstlisting}
	\caption{Код функций записи и получения значений окружения}
	\label{f:code_getsetEnv}
\end{figure}

\clearpage


Использование единого хранилища значений переменных для всей области видимости программы вносит некоторые ограничения.
Например в аргументах функции могут быть определены параметры, имена которых совпадают с объявленным ранее переменным.
Это некорректное поведение, так как первое объявленное значение будет перезаписано другим, переданным в функцию.
Пример кода такой ситуации приведен на рисунке~\ref{f:code_rewriteVarError}.

\begin{figure}[!htb]
	\centering
    \vspace{\toppaddingoffigure}
	\begin{lstlisting}[
        language=Go,
        xleftmargin=.08\textwidth,
        xrightmargin=.08\textwidth
    ]
x = 10;
f = func(x) {
    return x * 10;
}
f(5);
x; //5
\end{lstlisting}
	\caption{Пример кода некорректного поведения}
	\label{f:code_rewriteVarError}
\end{figure}

В данном коде, имя параметра функции совпадает с именем переменной -- x.
Переменная <<x>> объявлена со значением, равным 10.
Данное значение перезапишется значением, переданным в функцию при ее вызове, в данном случае значением, равным 5.
Таким образом, после вызова функции значение переменной <<x>> будет неявно изменено и станет равным 5.
Данная логика работы является ошибочной.

Решение лежит в выделении внутреннего окружения функции при её вызове.
Именно для такого случая в ранее рассмотренной структуре Env содержится ссылка на другой экземпляр структуры такого же типа.
Так, при вызове функции, необходимо создать новый экземпляр окружения, записать в него переданные аргументы и установить ссылку на внешнее окружение,
то из которого была вызвана функция. Такой подход позволит корректно выполнять вложенные функции и рекурсивные вызовы.

Для строк и массивов реализованы несколько встроенных функций:

\begin{itemize}
    \item len -- определение длины строки или массива;
    \item push -- добавление элемента в конец массива;
    \item first -- получение первого элемента массива;
    \item last -- получение последнего элемента массива.
\end{itemize}
% Код реализации встроенных функций приведен в приложении Г.

Хэш-карты реализованы на базе хэш-карт языка Go.
В качестве ключа могут быть следующие типы данных: строка, булево значение, число.
Так как данные типы представлены в виде внутренних объектов, нельзя брать тип Object в качестве ключа.
При занесении значения в хэш-карту и попытке последующего его получения ключи будут представлять разные экземпляры несмотря на одинаковое значение.
На рисунке~\ref{f:code_hashmapGetValueExample} приведен наглядный пример.

\begin{figure}[!htb]
	\centering
    \vspace{\toppaddingoffigure}
	\begin{lstlisting}[
        language=Go
    ]
X = { "name": "Bob" }
X["name"]     
\end{lstlisting}
	\caption{Пример получения значения хэш-карты}
	\label{f:code_hashmapGetValueExample}
\end{figure}

Строковой ключ «name» при попытке получить значение из карты не будет возвращать значение "Bob",
так как при вычислении выражения будет создан новый экземпляр объекта, представляющего строку.
Для решения этой проблемы можно использовать в качестве ключа строку, содержащую хэш от значения объекта.

Фрагмент кода объектной системы с реализацией хэш-карт представлен на рисунке~\ref{f:code_hashmapObject}.

Фрагмент кода функции вычисления целочисленного инфиксного выражения приведен на рисунке~\ref{f:code_evalIntInfixExpr}.

\clearpage

\begin{figure}[!htb]
	\centering
	\begin{lstlisting}[
        language=Go,
        xleftmargin=.08\textwidth,
        xrightmargin=.08\textwidth
    ]
type HashKey struct {
    Type  ObjectType
    Value uint64
}

type Hashable interface {
    HashKey() HashKey
}

type HashPair struct {
    Key   Object
    Value Object
}

type HashMap struct {
    Pairs map[HashKey]HashPair
}

func (h *HashMap) Type() ObjectType { return HASH_MAP_OBJ }
func (h *HashMap) ToString() string {
    var out bytes.Buffer

    pairs := []string{}
    for _, el := range h.Pairs {
        pairs = append(pairs, fmt.Sprintf("%s: %s", el.Key.ToString(), el.Value.ToString()))
    }

    out.WriteString("[")
    out.WriteString(strings.Join(pairs, ", "))
    out.WriteString("]")

    return out.String()
} 
\end{lstlisting}
	\caption{Фрагмент кода реализации хэш-карт}
	\label{f:code_hashmapObject}
\end{figure}

\clearpage

\begin{figure}[!htb]
	\centering
	\begin{lstlisting}[
        language=Go,
        xleftmargin=.08\textwidth,
        xrightmargin=.08\textwidth
    ]
func evalIntInfixExpr(op string, left object.Object, right object.Object) object.Object {
    lVal := left.(*object.Integer).Value
    rVal := right.(*object.Integer).Value
    switch op {
    case "+":
        return &object.Integer{Value: lVal + rVal}
    case "-":
        return &object.Integer{Value: lVal - rVal}
    case "*":
        return &object.Integer{Value: lVal * rVal}
    case "/":
        return &object.Integer{Value: lVal / rVal}
    case "%":
        return &object.Integer{Value: lVal % rVal}
    case "==":
        return boolToBooleanObj(lVal == rVal)
    case "!=":
        return boolToBooleanObj(lVal != rVal)
    case "<":
        return boolToBooleanObj(lVal < rVal)
    case ">":
        return boolToBooleanObj(lVal > rVal)
    case "<=":
        return boolToBooleanObj(lVal <= rVal)
    case ">=":
        return boolToBooleanObj(lVal >= rVal)
    default:
        return newError("unknown operator: %s %s %s", left.Type(), op, right.Type())
    }
}
\end{lstlisting}
	\caption{Фрагмент кода функции вычисления целочисленного инфиксного выражения}
	\label{f:code_evalIntInfixExpr}
\end{figure}

% Фрагменты кода исполнителя приведены в приложении В.

Пример успешного выполнения программы для указанных входных данных представлен на рисунке~\ref{f:evalSuccessExample}.

Входная строка: $5 + 1 * 20 / (5 + 5)$.

\clearpage

\begin{figure}[!htb]
	\centering
	\includegraphics[width=0.4\textwidth]{evaluator/evalSuccessExample.png}
	\caption{Пример успешного выполнения программы}
	\label{f:evalSuccessExample}
\end{figure}

Пример завершения работы программы с семантической ошибкой для указанных ниже входных данных показан на рисунке~\ref{f:evalErrorExample}.

Входная строка представлена на рисунке~\ref{f:code_evalErrorExample}.

\begin{figure}[!htb]
	\centering
    \vspace{\toppaddingoffigure}
	\begin{lstlisting}[
        language=Go,
        xleftmargin=.08\textwidth,
        xrightmargin=.08\textwidth
    ]
x = 5 + 1 * 20 / (5 + 5)
if (x > true) {
    return 1
}    
\end{lstlisting}
	\caption{Пример кода, содержащего ошибку}
	\label{f:code_evalErrorExample}
\end{figure}

\begin{figure}[!htb]
	\centering
	\includegraphics[width=0.8\textwidth]{evaluator/evalErrorExample.png}
	\caption{Пример завершения работы программы с ошибкой}
	\label{f:evalErrorExample}
\end{figure}

Листинг фрагментов кода реализации основных частей разработанного программного продукта представлен в приложении~\ref{ax:code}.